\chapter{Cryptographic Commitments and Proofs}

\section{Purpose}

Chapter 1 treated blockchains as replicated state machines. Chapter 2
identified assets and authority. Chapter 3 showed how transactions carry
inputs into execution. Chapter 4 studied keys, signatures, and authorization. Chapter 5 studies the cryptographic
mechanism that lets participants make compact claims about large bodies
of data:

\begin{quote}
A cryptographic commitment lets a system bind itself to data without
always revealing or transmitting all of that data.
\end{quote}

This is one of the quiet foundations of blockchain security.

A block hash commits to a block. A transaction hash commits to a
transaction. A Merkle root commits to a set or sequence of leaves. An
Ethereum \passthrough{\lstinline!state\_root!} commits to the post-block
world state. A \passthrough{\lstinline!receipts\_root!} commits to
receipts and logs. A bridge proof may claim that some event, storage
slot, or message was included in a source chain. A rollup output root
may commit to an L2 state. A data-availability commitment may claim that
data exists and can be sampled or reconstructed. A ZK verifier may check
a proof against public inputs that are themselves commitments.

These commitments are powerful because they let a verifier check a small
proof against a small root instead of downloading and trusting
everything.

But they are also dangerous when misunderstood.

A proof does not mean ``this is true'' in the abstract. It means:

\begin{lstlisting}
Given this commitment root,
and this verification algorithm,
and these cryptographic assumptions,
this claim is consistent with the data committed by that root.
\end{lstlisting}

That sentence should feel precise. Precision is the point.

\section{Learning Objectives}

By the end of this chapter, you should be able to:

\begin{itemize}
\tightlist
\item
  Explain preimage resistance, second-preimage resistance, and collision
  resistance in security-review terms.
\item
  Distinguish hashes, commitments, digests, identifiers, fingerprints,
  and proofs.
\item
  Explain why canonical encoding and domain separation are necessary for
  safe commitments.
\item
  Describe Merkle trees, Merkle Patricia tries, state roots, transaction
  roots, receipt roots, and inclusion proofs.
\item
  Explain what a proof actually proves and what it does not prove.
\item
  Identify proof-verification bugs involving wrong roots, stale roots,
  malformed proofs, replay, ambiguous encodings, and false trust
  assumptions.
\item
  Write a trust-minimized claim in the form: root, proof, verification
  algorithm, domain, finality, and security assumption.
\end{itemize}

\section{The Core Model}

A commitment system has three pieces:

\begin{lstlisting}
Data:
  The thing being committed to.

Commitment:
  A compact value derived from the data.

Opening or proof:
  Evidence that a specific claim is consistent with the commitment.
\end{lstlisting}

For a simple hash commitment:

\begin{lstlisting}
commitment = H(encoded_data)

proof = encoded_data

verify:
  H(encoded_data) == commitment
\end{lstlisting}

For a Merkle inclusion proof:

\begin{lstlisting}
commitment = merkle_root(all_leaves)

claim = leaf i has value v

proof = sibling hashes along the path from leaf i to root

verify:
  recompute root from v, i, and proof
  check recomputed_root == commitment
\end{lstlisting}

For a blockchain state proof:

\begin{lstlisting}
commitment = state_root in a block header

claim = account or storage key has value v

proof = trie nodes needed to verify that key/value path

verify:
  recompute path under the trie rules
  check resulting root == state_root
  check state_root is from an acceptable block
\end{lstlisting}

The proof never floats in empty space. It is always relative to a root
and a verification rule.


A minimal Merkle proof diagram should show the claim, the siblings, and the root:

\begin{lstlisting}
            root
           /    \
        h01      h23
       /  \      /  \
     h0   h1   h2   h3

claim:
  h2 is included

proof:
  sibling h3
  sibling h01

verify:
  h23   = H(h2 || h3)
  root' = H(h01 || h23)
  accept if root' == root
\end{lstlisting}

The security question is therefore:

\begin{lstlisting}
Which root is trusted?
Why is that root trusted?
What exact claim is proved against it?
Which algorithm verifies the proof?
Which assumptions would make the proof invalid, stale, or irrelevant?
\end{lstlisting}

\section{Hash Functions, Preimages, and Collision Resistance}

\subsection{A Cryptographic Hash Is Not Just a Checksum}

A cryptographic hash function maps arbitrary-length input to a
fixed-length output.

\begin{lstlisting}
digest = H(message)
\end{lstlisting}

NIST's Secure Hash Standard specifies SHA-family algorithms that
generate message digests, and NIST SP 800-107 Rev.~1, used here as
legacy terminology guidance, describes a hash function as mapping a bit
string of arbitrary length to a fixed-length bit string.\footnote{NIST,
  ``FIPS 180-4: Secure Hash Standard'',
  https://nvlpubs.nist.gov/nistpubs/FIPS/NIST.FIPS.180-4.pdf.}
\footnote{NIST, ``SP 800-107 Revision 1: Recommendation for Applications
  Using Approved Hash Algorithms'', legacy guidance whose supplementary
  material has moved to NIST hash-function resources,
  https://nvlpubs.nist.gov/nistpubs/Legacy/SP/nistspecialpublication800-107r1.pdf.}

For blockchain security, the important practical properties are:

\begin{itemize}
\tightlist
\item
  Determinism: the same input gives the same output.
\item
  Preimage resistance: given a digest, it should be infeasible to find
  an input that hashes to it.
\item
  Second-preimage resistance: given one input, it should be infeasible
  to find a different input with the same digest.
\item
  Collision resistance: it should be infeasible to find any two
  different inputs with the same digest.
\item
  Efficient verification: anyone can recompute and compare.
\end{itemize}

These properties let a small digest act as a compact fingerprint for
much larger data.

But a cryptographic hash is not encryption. It does not hide low-entropy
data. It does not authenticate who produced the data. It does not prove
that the data is meaningful, safe, fresh, or final.

It proves only a relationship:

\begin{lstlisting}
This input hashes to this output.
\end{lstlisting}

Everything else comes from context.

\subsection{Preimage Resistance}

Preimage resistance means that given a hash output, an attacker should
not be able to find an input that produces it.

\begin{lstlisting}
Given y = H(x),
it should be infeasible to find any x' such that H(x') == y.
\end{lstlisting}

NIST SP 800-107 Rev.~1 describes preimage resistance as the property
that, given a randomly chosen hash value, finding an input that maps to
it is computationally infeasible. For an L-bit hash function, expected
preimage strength is generally L bits.\footnote{NIST, ``SP 800-107
  Revision 1: Recommendation for Applications Using Approved Hash
  Algorithms'', legacy guidance whose supplementary material has moved
  to NIST hash-function resources,
  https://nvlpubs.nist.gov/nistpubs/Legacy/SP/nistspecialpublication800-107r1.pdf.}

In blockchain systems, preimage resistance appears in:

\begin{itemize}
\tightlist
\item
  Hashlocks.
\item
  Commit-reveal schemes.
\item
  Address derivation.
\item
  Transaction and block identifiers.
\item
  Content-addressed storage.
\item
  Proof-of-work puzzles.
\item
  Merkle roots and trie node references.
\end{itemize}

Security review question:

\begin{lstlisting}
Is the committed input unpredictable enough that attackers cannot brute-force it?
\end{lstlisting}

This question matters because preimage resistance of the hash function
is not enough if the input space is tiny.

Example:

\begin{lstlisting}
commitment = H("yes")
\end{lstlisting}

This hides nothing. An attacker can hash \passthrough{\lstinline!"yes"!}
and \passthrough{\lstinline!"no"!} and compare. A hash commitment to a
low-entropy value needs a secret salt or nonce if hiding matters.

Better:

\begin{lstlisting}
commitment = H("vote-v1" || voter || proposal_id || vote || random_256_bit_salt)
\end{lstlisting}

The hash function may be strong, but the protocol must give it enough
entropy and context.

\subsection{Second-Preimage Resistance}

Second-preimage resistance means that after seeing a specific input, an
attacker should not be able to find a different input with the same
digest.

\begin{lstlisting}
Given x,
it should be infeasible to find x' != x such that H(x') == H(x).
\end{lstlisting}

This property matters when a commitment has already been made to a known
object.

Examples:

\begin{itemize}
\tightlist
\item
  Replacing a transaction with a different transaction that has the same
  transaction ID.
\item
  Replacing a block body while preserving a block commitment.
\item
  Replacing a contract artifact while preserving an expected source
  hash.
\item
  Replacing a leaf value while preserving a Merkle root.
\end{itemize}

If second-preimage resistance fails, a commitment no longer binds the
committer to the original data.

Security review question:

\begin{lstlisting}
Can the prover substitute a different object while preserving the committed digest?
\end{lstlisting}

Usually the answer should be ``computationally infeasible.'' But that
answer depends on the correct hash function, correct encoding, correct
domain separation, and no exploitable ambiguity in what is being hashed.

\subsection{Collision Resistance}

Collision resistance means that an attacker should not be able to find
any two different inputs with the same digest.

\begin{lstlisting}
Find x != x' such that H(x) == H(x').
\end{lstlisting}

NIST SP 800-107 Rev.~1 states that for an L-bit hash function, the
expected collision-resistance strength is L/2 bits because of
birthday-bound attacks; for SHA-256, this corresponds to an expected
128-bit collision strength.\footnote{NIST, ``SP 800-107 Revision 1:
  Recommendation for Applications Using Approved Hash Algorithms'',
  legacy guidance whose supplementary material has moved to NIST
  hash-function resources,
  https://nvlpubs.nist.gov/nistpubs/Legacy/SP/nistspecialpublication800-107r1.pdf.}

Collision resistance matters when attackers can choose both inputs.

Examples:

\begin{itemize}
\tightlist
\item
  A malicious builder creates two different payloads with the same
  commitment.
\item
  A contract accepts one document during review and a different document
  during execution.
\item
  Two typed messages collide because encoding is ambiguous.
\item
  A Merkle tree allows a leaf to be interpreted as an internal node.
\end{itemize}

Collision attacks are not the only hash risk, but they are central to
binding commitments. If an attacker can produce two valid openings for
one commitment, the commitment is not binding.

\subsection{Hash Security Depends on the Application}

The same hash function may be used for different security goals.

{\def\LTcaptype{none} % do not increment counter
\begin{longtable}[]{@{}
  >{\raggedright\arraybackslash}p{(\linewidth - 4\tabcolsep) * \real{0.3333}}
  >{\raggedright\arraybackslash}p{(\linewidth - 4\tabcolsep) * \real{0.3333}}
  >{\raggedright\arraybackslash}p{(\linewidth - 4\tabcolsep) * \real{0.3333}}@{}}
\toprule\noalign{}
\begin{minipage}[b]{\linewidth}\raggedright
Use
\end{minipage} & \begin{minipage}[b]{\linewidth}\raggedright
Needed property
\end{minipage} & \begin{minipage}[b]{\linewidth}\raggedright
Failure mode
\end{minipage} \\
\midrule\noalign{}
\endhead
\bottomrule\noalign{}
\endlastfoot
Hashlock & Preimage resistance & Attacker claims without knowing
secret \\
File fingerprint & Second-preimage resistance & Attacker swaps file for
another with same hash \\
Digital signature prehash & Collision and second-preimage resistance &
Attacker signs one message, substitutes another \\
Merkle root & Collision and second-preimage resistance & Attacker proves
false membership \\
Commit-reveal & Binding and possibly hiding & Attacker changes value or
brute-forces committed value \\
Content addressing & Second-preimage resistance & Wrong content served
under expected address \\
Proof-of-work & Preimage-style search difficulty under target & Attacker
finds cheap block headers below target \\
\end{longtable}
}

This is why ``we hash it'' is not a security argument. You must say
which property the protocol needs and why the construction gives it.

\subsection{Hashes Are Identifiers, Not Authority}

Hashes often act as identifiers:

\begin{itemize}
\tightlist
\item
  Transaction hash.
\item
  Block hash.
\item
  Code hash.
\item
  Storage root.
\item
  Content hash.
\item
  Proposal hash.
\item
  Message hash.
\end{itemize}

But an identifier does not by itself authorize anything. A transaction
hash identifies a transaction. It does not prove that the transaction is
final, successful, economically safe, or submitted by the correct user
unless the verifier checks the surrounding evidence.

Bad:

\begin{lstlisting}
The transaction hash proves the transfer happened.
\end{lstlisting}

Better:

\begin{lstlisting}
The transaction hash identifies a transaction. To prove the transfer happened, verify that the transaction is included in the canonical chain, executed successfully, emitted or caused the expected state change, and is final enough for the relying system.
\end{lstlisting}

Hashes are pointers into evidence. They are not the evidence by
themselves.

\subsection{Hash Functions Do Not Fix Bad Encoding}

Hash functions process bytes. They do not understand your data model.

If two different meanings serialize to the same bytes, the hash cannot
save you.

Example of ambiguous concatenation:

\begin{lstlisting}
H("ab" || "c") == H("a" || "bc")
\end{lstlisting}

Both hash the same bytes:

\begin{lstlisting}
"abc"
\end{lstlisting}

Safe encodings include:

\begin{itemize}
\tightlist
\item
  Length prefixes.
\item
  Type tags.
\item
  Version tags.
\item
  Fixed-width fields where appropriate.
\item
  Canonical serialization.
\item
  Domain separators.
\item
  Explicit chain and contract identifiers.
\end{itemize}

The security property is not:

\begin{lstlisting}
H(data) is collision resistant.
\end{lstlisting}

It is:

\begin{lstlisting}
H(canonical_encode(domain, type, version, fields)) binds to exactly one meaning.
\end{lstlisting}

Many proof and signature bugs are encoding bugs wearing cryptographic
clothing.

\section{Commitments, Binding Assumptions, and Domain Separation}

\subsection{A Commitment Binds a Claim to Data}

A commitment is a compact value that binds the committer to data.

In cryptography, commitment schemes are usually discussed through two
properties:

\begin{itemize}
\tightlist
\item
  Binding: after committing, the committer cannot open the commitment to
  a different value.
\item
  Hiding: before opening, the verifier cannot learn the committed value.
\end{itemize}

Blockchain commitments often care primarily about binding and
verifiability. Hiding may or may not be required.

Examples:

{\def\LTcaptype{none} % do not increment counter
\begin{longtable}[]{@{}
  >{\raggedright\arraybackslash}p{(\linewidth - 6\tabcolsep) * \real{0.2500}}
  >{\raggedright\arraybackslash}p{(\linewidth - 6\tabcolsep) * \real{0.2500}}
  >{\raggedright\arraybackslash}p{(\linewidth - 6\tabcolsep) * \real{0.2500}}
  >{\raggedright\arraybackslash}p{(\linewidth - 6\tabcolsep) * \real{0.2500}}@{}}
\toprule\noalign{}
\begin{minipage}[b]{\linewidth}\raggedright
Commitment
\end{minipage} & \begin{minipage}[b]{\linewidth}\raggedright
Usually hiding?
\end{minipage} & \begin{minipage}[b]{\linewidth}\raggedright
Usually binding?
\end{minipage} & \begin{minipage}[b]{\linewidth}\raggedright
Opening/proof
\end{minipage} \\
\midrule\noalign{}
\endhead
\bottomrule\noalign{}
\endlastfoot
\passthrough{\lstinline!H(transaction)!} & No & Yes, assuming hash
security and canonical encoding & Transaction bytes \\
Block hash & No & Yes & Block header/body validation \\
Merkle root & No & Yes & Merkle path \\
State root & No & Yes & Trie proof and state-transition validation \\
Commit-reveal vote hash & Should be & Should be & Vote plus salt \\
KZG commitment & Not necessarily & Yes under algebraic assumptions &
Evaluation proof \\
\end{longtable}
}

The important lesson:

\begin{lstlisting}
Commitment does not automatically mean secrecy.
\end{lstlisting}

A Merkle root commits to leaves. It usually does not hide them if the
leaves are public or guessable. A hash of a small value is easy to
brute-force. A public transaction hash does not hide the transaction
once the transaction is known.

\subsection{Binding Depends on the Whole Construction}

Binding is not supplied by the hash function alone. It depends on:

\begin{itemize}
\tightlist
\item
  The cryptographic primitive.
\item
  The exact bytes committed.
\item
  Canonical serialization.
\item
  Domain separation.
\item
  Tree construction rules.
\item
  Verification rules.
\item
  The root source.
\item
  Finality assumptions.
\end{itemize}

Suppose a bridge accepts:

\begin{lstlisting}
message_hash = H(amount || recipient || nonce)
\end{lstlisting}

Questions:

\begin{itemize}
\tightlist
\item
  Is \passthrough{\lstinline!amount!} fixed-width?
\item
  Is \passthrough{\lstinline!recipient!} fixed-width?
\item
  Is \passthrough{\lstinline!nonce!} fixed-width?
\item
  Is the source chain included?
\item
  Is the destination chain included?
\item
  Is the bridge contract included?
\item
  Is the message type included?
\item
  Is the token address included?
\item
  Is the version included?
\end{itemize}

If not, the hash may be binding bytes but not binding the intended
meaning.

The better construction is closer to:

\begin{lstlisting}
message_hash = H(
  "bridge-message-v1" ||
  source_chain_id ||
  source_bridge_address ||
  destination_chain_id ||
  destination_bridge_address ||
  token_address ||
  recipient ||
  amount ||
  nonce
)
\end{lstlisting}

This is not merely cleaner. It prevents cross-protocol and cross-domain
confusion.

\subsection{Domain Separation Prevents Context Confusion}

Domain separation means that commitments for one purpose cannot be
validly interpreted as commitments for another purpose.

At minimum, domain separation asks:

\begin{lstlisting}
What protocol is this for?
What chain or domain is this for?
What contract, program, or module is this for?
What version is this for?
What type of object is this?
What operation is this authorizing or proving?
\end{lstlisting}

Without domain separation, the same hash or signature may be reusable in
another context.

Examples:

\begin{itemize}
\tightlist
\item
  A message signed for chain A is accepted on chain B.
\item
  A withdrawal proof for token X is accepted for token Y.
\item
  A vote commitment is interpreted as an order commitment.
\item
  A Merkle leaf in one tree type is interpreted as an internal node in
  another.
\item
  A bridge message for one destination contract is replayed into
  another.
\item
  A testnet signature works on mainnet.
\end{itemize}

Domain separation is not optional ceremony. It is how a protocol says:

\begin{lstlisting}
This commitment means this one thing here, and not something else somewhere else.
\end{lstlisting}

\subsection{Canonical Encoding Is Part of the Commitment}

If two honest implementations encode the same object differently, they
will compute different commitments. If two different objects encode to
the same bytes, the commitment is ambiguous.

Canonical encoding requires:

\begin{itemize}
\tightlist
\item
  One field order.
\item
  One integer representation.
\item
  One byte-order convention.
\item
  One string normalization policy, or no strings where bytes are
  expected.
\item
  One rule for optional fields.
\item
  One rule for empty values.
\item
  One rule for lists, maps, and sorting.
\item
  One versioning scheme.
\end{itemize}

Ethereum uses Recursive Length Prefix encoding in its execution-layer
trie construction, and the Ethereum consensus specs define Simple
Serialize for consensus-layer objects and merkleization.\footnote{ethereum.org,
  ``Merkle Patricia Trie'',
  https://ethereum.org/developers/docs/data-structures-and-encoding/patricia-merkle-trie/.}
\footnote{Ethereum Consensus Specs, ``SimpleSerialize (SSZ)'',
  https://ethereum.github.io/consensus-specs/ssz/simple-serialize/.}

Security consequences:

\begin{itemize}
\tightlist
\item
  A proof verifier must decode exactly according to the committed
  format.
\item
  A bridge verifier must not accept ``equivalent'' but noncanonical
  encodings.
\item
  A signature verifier must bind the same bytes the user reviewed.
\item
  A client implementation must not disagree with other clients about
  serialization.
\end{itemize}

If serialization is ambiguous, the cryptography is standing on wet
ground.

\subsection{Hiding Requires Entropy}

A hash commitment hides only if the committed value is hard to guess.

Unsafe:

\begin{lstlisting}
commitment = H(bid_amount)
\end{lstlisting}

If bid amounts are predictable, every bidder can brute-force likely
values.

Better:

\begin{lstlisting}
commitment = H("sealed-bid-v1" || auction_id || bidder || bid_amount || salt)
\end{lstlisting}

The salt must be unpredictable and large enough. A four-digit salt is
not enough. A 256-bit random salt is usually the kind of magnitude
expected for high-value commitments.

Hiding analysis asks:

\begin{lstlisting}
Can an attacker enumerate possible committed values?
Is there a random salt?
Is the salt secret until reveal?
Is the salt unique?
Is the commitment bound to this auction, bidder, chain, and contract?
Can revealing in one context leak or authorize in another?
\end{lstlisting}

Commit-reveal protocols fail surprisingly often because the commit phase
binds too little or hides too little.

\subsection{Commitments Can Become Stale}

A commitment may be correct and still unsafe to rely on.

Examples:

\begin{itemize}
\tightlist
\item
  A state root from an old block proves an old balance, not the current
  balance.
\item
  A transaction inclusion proof from an unfinalized block may disappear
  in a reorg.
\item
  A bridge proof may be valid for a source-chain event that is not final
  enough.
\item
  A rollup output root may be disputed during a challenge window.
\item
  A price commitment may be stale even if correctly signed.
\item
  A proof from a pre-upgrade contract may not mean the same thing after
  upgrade.
\end{itemize}

Every commitment has a time and domain.

Security review asks:

\begin{lstlisting}
At what block, slot, epoch, version, or state root is this claim true?
Is that root final enough?
Is the relying protocol allowed to use old data?
Does the commitment expire?
Can a fresher contradictory commitment exist?
\end{lstlisting}

Freshness is not a cryptographic property. It is a protocol property.

\section{Merkle Trees, Patricia Tries, State Roots, and Inclusion Proofs}

\subsection{Merkle Trees Compress Many Claims Into One Root}

A Merkle tree hashes leaves into parent nodes until one root remains.

Simple binary example:

\begin{lstlisting}
L0 = H("leaf" || data0)
L1 = H("leaf" || data1)
L2 = H("leaf" || data2)
L3 = H("leaf" || data3)

P0 = H("node" || L0 || L1)
P1 = H("node" || L2 || L3)

root = H("node" || P0 || P1)
\end{lstlisting}

The root commits to all leaves and their positions. A proof for
\passthrough{\lstinline!data2!} needs:

\begin{lstlisting}
data2
L3
P0
position information: right child of P1, P1 right child of root
\end{lstlisting}

The verifier recomputes:

\begin{lstlisting}
L2 = H("leaf" || data2)
P1 = H("node" || L2 || L3)
root' = H("node" || P0 || P1)
accept if root' == root
\end{lstlisting}

RFC 6962, the Certificate Transparency specification, uses Merkle hash
trees to let clients verify that data is included in a log and to check
consistency between tree sizes.\footnote{IETF, ``RFC 6962: Certificate
  Transparency'', https://datatracker.ietf.org/doc/html/rfc6962.}
Bitcoin uses Merkle trees so a lightweight client can verify that a
transaction is included in a block without downloading the full
block.\footnote{Satoshi Nakamoto, ``Bitcoin: A Peer-to-Peer Electronic
  Cash System'', https://bitcoin.org/bitcoin.pdf.}

For security analysis, the important feature is:

\begin{lstlisting}
A small proof can connect one leaf to one trusted root.
\end{lstlisting}

\subsection{Inclusion Proofs Prove Membership Relative to a Root}

An inclusion proof proves:

\begin{lstlisting}
This leaf is included in the committed data structure represented by this root.
\end{lstlisting}

It does not prove:

\begin{itemize}
\tightlist
\item
  The root is canonical.
\item
  The root is final.
\item
  The leaf is semantically valid.
\item
  The leaf has not been superseded.
\item
  The transaction succeeded.
\item
  The event means what the verifier thinks it means.
\item
  The same claim has not already been consumed elsewhere.
\end{itemize}

Those facts require additional checks.

For a Bitcoin SPV-style payment check, a Merkle branch can show
transaction inclusion in a block. But the verifier still needs the block
header chain and enough confirmation confidence. The Bitcoin
whitepaper's simplified payment verification section relies on block
headers and proof-of-work chain evidence, not a Merkle proof
alone.\footnote{Satoshi Nakamoto, ``Bitcoin: A Peer-to-Peer Electronic
  Cash System'', https://bitcoin.org/bitcoin.pdf.}

For a bridge, an event proof can show that an event log was included in
a source-chain receipt trie. But the bridge still needs:

\begin{itemize}
\tightlist
\item
  Correct source chain.
\item
  Correct source contract.
\item
  Correct event signature and fields.
\item
  Correct block header.
\item
  Correct finality threshold.
\item
  Correct replay protection.
\item
  Correct destination-chain message handling.
\end{itemize}

The Merkle proof is one piece.

\subsection{Tree Shape and Leaf Meaning Matter}

Merkle verification must match the exact tree construction.

Important questions:

\begin{itemize}
\tightlist
\item
  Are leaves ordered?
\item
  Is the tree binary, sparse, Patricia, or another structure?
\item
  Are internal nodes domain-separated from leaves?
\item
  Are leaf values hashed before insertion?
\item
  Are duplicate leaves allowed?
\item
  What happens with an odd number of leaves?
\item
  Does the proof include positions, or does the verifier sort pairs?
\item
  Does membership include an index?
\item
  Are absence proofs supported?
\item
  Is the proof for a set, sequence, map, or key/value store?
\end{itemize}

A proof for ``leaf value X exists somewhere'' is weaker than a proof for
``leaf value X exists at index i.'' A proof for ``key k maps to value
v'' is different from a proof for ``value v appears somewhere.''

Security failure example:

\begin{lstlisting}
Verifier accepts sorted pair hashing.
Protocol intended ordered leaves.
Attacker reorders siblings and proves a different position.
\end{lstlisting}

Another:

\begin{lstlisting}
Verifier hashes raw leaf bytes.
Internal node bytes can be shaped like leaf bytes.
Attacker exploits leaf/internal ambiguity.
\end{lstlisting}

Domain tags such as \passthrough{\lstinline!"leaf"!} and
\passthrough{\lstinline!"node"!} are a simple defense in many
constructions. More generally, the committed encoding must make each
node type unambiguous.

\subsection{Patricia Tries Commit to Key/Value State}

A Merkle tree is good for proving that a leaf appears in a sequence or
set. A blockchain state often needs a map:

\begin{lstlisting}
key -> value
\end{lstlisting}

Ethereum uses a modified Merkle Patricia Trie for execution-layer state.
Ethereum.org describes Ethereum's state as encoded into a special
Merkle-tree-like structure, where a root value can prove facts about the
data, and explains that Ethereum's structure is a modified
Merkle-Patricia trie designed for efficient state retrieval.\footnote{ethereum.org,
  ``Merkle Patricia Trie'',
  https://ethereum.org/developers/docs/data-structures-and-encoding/patricia-merkle-trie/.}

At a high level:

\begin{lstlisting}
state root
  commits to account trie

account trie entry
  commits to account fields, including storage root

storage root
  commits to contract storage trie
\end{lstlisting}

Ethereum blocks contain commitments such as:

\begin{itemize}
\tightlist
\item
  \passthrough{\lstinline!stateRoot!}: root of the final state after
  applying transactions.
\item
  \passthrough{\lstinline!transactionsRoot!}: root of the transaction
  trie.
\item
  \passthrough{\lstinline!receiptsRoot!}: root of the receipts trie.
\end{itemize}

Ethereum.org's block documentation lists these roots in block and
execution-payload structures, and notes that clients re-execute
transactions to ensure the new state matches the block's state root; the
Yellow Paper formalizes the block/state-transition framing and trie-root
fields.\footnote{ethereum.org, ``Blocks'',
  https://ethereum.org/developers/docs/blocks/.} \footnote{Gavin Wood,
  ``Ethereum: A Secure Decentralised Generalised Transaction Ledger'',
  https://ethereum.github.io/yellowpaper/paper.pdf.}

This gives us the state-machine connection:

\begin{lstlisting}
parent_state_root
+ ordered transactions
+ execution rules
= new_state_root
\end{lstlisting}

If a block claims a \passthrough{\lstinline!stateRoot!}, honest
validators re-execute the transactions and check that the computed root
matches. If they disagree about execution rules or trie construction,
consensus can split.

\subsection{State Proofs and Storage Proofs}

A state proof proves that an account or storage value is included in a
committed state root.

EIP-1186 specifies the \passthrough{\lstinline!eth\_getProof!} RPC
method for retrieving Merkle proofs for an account and storage
slots.\footnote{Ethereum Improvement Proposals, ``EIP-1186: RPC-Method
  to get Merkle Proofs - eth\_getProof'',
  https://eips.ethereum.org/EIPS/eip-1186.} The method returns account
fields and proof nodes that allow a verifier to check claims against
Ethereum's state trie.

A simplified storage-proof flow:

\begin{lstlisting}
Inputs:
  trusted block header
  stateRoot from header
  account address
  storage slot key
  claimed storage value
  account proof
  storage proof

Verify account proof:
  account address -> account object under stateRoot

Extract storageRoot from account object

Verify storage proof:
  storage slot key -> claimed value under storageRoot

Accept if:
  both trie proofs verify
  block header is trusted/final enough
  chain/domain is correct
  claimed value has the expected meaning
\end{lstlisting}

The proof does not prove that the storage value is current forever. It
proves the value at that state root.

This distinction matters for:

\begin{itemize}
\tightlist
\item
  Bridges.
\item
  Cross-chain governance.
\item
  Historical queries.
\item
  On-chain light clients.
\item
  Rollup withdrawals.
\item
  Oracle attestations.
\item
  Proof-of-reserves systems.
\end{itemize}

State proofs are often the difference between trusting an RPC response
and verifying a claim against a committed root.

\subsection{Receipt Proofs and Event Proofs}

Events live in transaction logs, and logs are represented in receipts.
Ethereum block headers include a \passthrough{\lstinline!receiptsRoot!},
committing to the receipts trie.\footnote{ethereum.org, ``Blocks'',
  https://ethereum.org/developers/docs/blocks/.}

An event proof usually needs to show:

\begin{lstlisting}
block header commits to receiptsRoot
receiptsRoot commits to transaction receipt
receipt contains log
log was emitted by expected contract
log topics and data match expected event
block is canonical and final enough
event has not already been consumed
\end{lstlisting}

Common mistakes:

\begin{itemize}
\tightlist
\item
  Verifying the event data but not the emitting contract.
\item
  Verifying the event but not the source chain.
\item
  Verifying inclusion but not finality.
\item
  Trusting a relayer's JSON instead of verifying the receipt proof.
\item
  Forgetting replay protection on the destination chain.
\item
  Treating event existence as proof that the corresponding application
  invariant is true.
\end{itemize}

An event is a statement emitted by code. It is not automatically the
state itself. Chapter 3's rule still holds:

\begin{lstlisting}
Events are evidence.
State is authority.
\end{lstlisting}

\subsection{SSZ Merkleization and Consensus Objects}

Ethereum's consensus layer uses Simple Serialize (SSZ), which defines
serialization, merkleization, and generalized indices for consensus
objects.\footnote{Ethereum Consensus Specs, ``SimpleSerialize (SSZ)'',
  https://ethereum.github.io/consensus-specs/ssz/simple-serialize/.}

This matters because proof systems are not interchangeable.

An Ethereum execution-layer Merkle Patricia proof and an Ethereum
consensus-layer SSZ proof have different:

\begin{itemize}
\tightlist
\item
  Serialization.
\item
  Tree structure.
\item
  Key/indexing model.
\item
  Hashing rules.
\item
  Object schemas.
\item
  Verification algorithms.
\end{itemize}

Security review question:

\begin{lstlisting}
Which commitment scheme is this proof for?
\end{lstlisting}

If the answer is vague, proof verification is already in trouble.

\section{Proof Verification, Proof Forgery Risk, and Trust-Minimized Claims}

\subsection{A Proof Is Only as Good as Its Root}

Every proof is relative to a root.

Bad verification:

\begin{lstlisting}
verify(proof, claim) == true
\end{lstlisting}

Better verification:

\begin{lstlisting}
verify(root, proof, claim, domain, rules) == true
\end{lstlisting}

Best verification also asks:

\begin{lstlisting}
Where did root come from?
Is root canonical?
Is root final enough?
Is root from the expected chain, block, contract, epoch, version, and fork?
Is the verifying code using the exact committed encoding?
Has this claim already been consumed?
\end{lstlisting}

A valid proof against an untrusted root is meaningless. An invalid root
can commit to any attacker's preferred universe.

\subsection{Proof Verification Is a State Transition}

When a smart contract, bridge, rollup, wallet, or backend verifies a
proof and acts on it, proof verification becomes part of the state
machine.

Example:

\begin{lstlisting}
claimWithdrawal(proof, message):
  require(verifyProof(trustedRoot, proof, message))
  require(!consumed[message.id])
  consumed[message.id] = true
  transfer(message.amount, message.recipient)
\end{lstlisting}

The proof verifier is not passive. It gates value movement.

Security review must inspect:

\begin{itemize}
\tightlist
\item
  The root registry.
\item
  The proof parser.
\item
  The encoding rules.
\item
  The hash function.
\item
  The domain separator.
\item
  The finality rule.
\item
  The consumed-message state.
\item
  The action triggered after verification.
\item
  The failure behavior if verification reverts or runs out of gas.
\end{itemize}

Proof verification bugs are often bridge bugs, withdrawal bugs,
governance bugs, and oracle bugs.

\subsection{Common Proof Verification Failures}

{\def\LTcaptype{none} % do not increment counter
\begin{longtable}[]{@{}
  >{\raggedright\arraybackslash}p{(\linewidth - 4\tabcolsep) * \real{0.3333}}
  >{\raggedright\arraybackslash}p{(\linewidth - 4\tabcolsep) * \real{0.3333}}
  >{\raggedright\arraybackslash}p{(\linewidth - 4\tabcolsep) * \real{0.3333}}@{}}
\toprule\noalign{}
\begin{minipage}[b]{\linewidth}\raggedright
Failure
\end{minipage} & \begin{minipage}[b]{\linewidth}\raggedright
Description
\end{minipage} & \begin{minipage}[b]{\linewidth}\raggedright
Result
\end{minipage} \\
\midrule\noalign{}
\endhead
\bottomrule\noalign{}
\endlastfoot
Wrong root & Proof verifies against attacker-supplied or stale root &
False claims accepted \\
Wrong domain & Proof from chain A accepted for chain B & Cross-domain
replay \\
Wrong contract & Event from fake contract accepted & Forged bridge or
oracle event \\
Wrong encoding & Verifier hashes bytes differently from producer & False
rejection or false acceptance \\
Missing position & Leaf value proved without index/key & Reordered or
duplicate claim accepted \\
Missing nonmembership semantics & Absence proof interpreted incorrectly
& Claims allowed despite existing state \\
Missing replay protection & Same proof consumed twice & Double mint or
double withdrawal \\
Insufficient finality & Proof from reorged block accepted & Unbacked
release \\
Malformed proof parser & Invalid nodes accepted or decoder confused &
Proof forgery or denial of service \\
Gas/resource failure & Valid proofs too costly to verify & Liveness
failure \\
Semantic mismatch & Proof is valid but proves the wrong thing & Unsafe
state transition \\
\end{longtable}
}

The last row is especially common. A proof can be mathematically valid
and still not support the action the protocol takes.

\subsection{Proof of Inclusion Is Not Proof of Validity}

Suppose a bridge sees a source-chain event:

\begin{lstlisting}
Deposit(user, amount, nonce)
\end{lstlisting}

An inclusion proof can show that this event was in a receipt committed
by a source-chain block.

It does not automatically prove:

\begin{itemize}
\tightlist
\item
  The deposit contract was not upgradeable to malicious code.
\item
  The deposited token was genuine.
\item
  The deposit was not later reverted by a reorg.
\item
  The source chain had sufficient finality.
\item
  The event fields were interpreted correctly.
\item
  The message was intended for this destination.
\item
  The user has not already claimed on the destination chain.
\end{itemize}

The proof answers an inclusion question. The protocol's security
requires additional validity, finality, and replay checks.

This distinction is one of the most important in cross-chain security:

\begin{lstlisting}
Inclusion answers "was this committed?"
Validity answers "should this committed thing be trusted for this action?"
\end{lstlisting}

\subsection{Trust-Minimized Does Not Mean Trust-Free}

Trust-minimized systems reduce the number or strength of assumptions a
user must accept. They do not eliminate assumptions.

A storage proof can remove the need to trust an RPC provider's answer
about a storage slot. But the verifier still trusts:

\begin{itemize}
\tightlist
\item
  The block header root.
\item
  The consensus/finality mechanism that accepted that block.
\item
  The hash and trie assumptions.
\item
  The proof verification implementation.
\item
  The meaning of the storage slot.
\item
  The freshness requirement for the relying action.
\end{itemize}

A light client can reduce reliance on a full node. It still trusts a
sync committee, validator quorum, proof-of-work chain, or other
consensus evidence depending on design.

A ZK proof can reduce the need to re-execute computation. It still
trusts the statement, circuit, verifier, public inputs, setup
assumptions if any, and integration logic.

Serious phrasing:

\begin{lstlisting}
This system verifies storage membership against finalized Ethereum headers,
so it does not trust an RPC provider for the storage value.
It still assumes Ethereum finality, correct header verification,
correct MPT verification, and correct interpretation of the storage slot.
\end{lstlisting}

That is how to write trust-minimized claims without fantasy.

\subsection{Proofs Need Freshness and Expiration}

Proofs often age badly.

Examples:

\begin{itemize}
\tightlist
\item
  A proof that an account had a balance at block N does not prove it has
  that balance at block N + 100.
\item
  A proof that a message was unconsumed at block N does not prove it is
  unconsumed now.
\item
  A proof that a validator set was active in epoch E may not apply in
  epoch E + 1.
\item
  A proof that an oracle price was fresh at one block may be stale by
  execution time.
\item
  A proof generated before an upgrade may not correspond to post-upgrade
  semantics.
\end{itemize}

Defenses:

\begin{itemize}
\tightlist
\item
  Include block number, slot, epoch, or timestamp in the verified claim.
\item
  Enforce maximum age.
\item
  Enforce finality minimums.
\item
  Store consumed nullifiers or message IDs.
\item
  Bind proofs to destination chain, contract, and version.
\item
  Require expected state checks at execution time.
\item
  Make stale proof failure explicit and safe.
\end{itemize}

Freshness belongs in the verifier, not in off-chain hope.

\subsection{Proof Verification Can Be a Denial-of-Service Surface}

Proof verification may be expensive.

Attackers may submit:

\begin{itemize}
\tightlist
\item
  Very large proofs.
\item
  Deep trie paths.
\item
  Malformed encodings.
\item
  Duplicate nodes.
\item
  Proofs that almost verify but fail late.
\item
  Proofs that trigger worst-case parser behavior.
\item
  Proofs for stale roots that waste computation.
\end{itemize}

On-chain proof verification must consider:

\begin{itemize}
\tightlist
\item
  Gas or compute limits.
\item
  Maximum proof size.
\item
  Maximum path length.
\item
  Canonical encoding checks.
\item
  Early rejection.
\item
  Memory allocation.
\item
  Hash cost.
\item
  Reentrancy or callbacks after proof acceptance.
\item
  Whether failed proofs cost the attacker enough.
\end{itemize}

A verifier that is cryptographically correct but economically unusable
can still fail security goals through liveness.

\subsection{The Verification Checklist}

For any proof-based protocol, write the verification checklist
explicitly:

\begin{lstlisting}
Claim:
  What exact fact is being proved?

Root:
  Which commitment root is used?

Root source:
  Where does the root come from?

Finality:
  Why is the root canonical and final enough?

Domain:
  Which chain, contract, program, module, version, and operation is this proof for?

Encoding:
  Which canonical serialization and hashing rules apply?

Proof type:
  Inclusion, absence, consistency, state, receipt, storage, light-client, ZK, or other?

Verifier:
  Which algorithm and implementation checks it?

Freshness:
  How old can the root or proof be?

Replay:
  What state marks the claim consumed?

Action:
  What state transition happens after verification?

Failure:
  What happens if proof verification fails or is too expensive?
\end{lstlisting}

If a protocol cannot answer these, the proof system is not yet a
security system. It is a cryptographic decoration.

\section{Worked Example: Verifying a Cross-Chain Deposit Claim}

Suppose a destination-chain bridge wants to mint wrapped tokens after a
user deposits tokens into a source-chain bridge.

Naive design:

\begin{lstlisting}
Relayer submits:
  user
  amount
  source_tx_hash

Destination bridge:
  if relayer is trusted:
    mint(user, amount)
\end{lstlisting}

This trusts the relayer for safety.

More trust-minimized design:

\begin{lstlisting}
Relayer submits:
  source block header
  receipt proof
  deposit event fields
  transaction index / receipt index
  log index

Destination bridge verifies:
  source block header is accepted and final enough
  receiptsRoot is in the block header
  receipt proof verifies against receiptsRoot
  receipt contains Deposit event
  event was emitted by expected source bridge
  event includes destination chain and destination bridge
  event includes token, recipient, amount, nonce
  nonce has not been consumed
\end{lstlisting}

Then:

\begin{lstlisting}
consumed[source_chain_id][source_bridge][nonce] = true
mint(recipient, amount)
\end{lstlisting}

Now compare the trust assumptions.

Naive design:

\begin{lstlisting}
The relayer tells the truth.
\end{lstlisting}

Proof-based design:

\begin{lstlisting}
The relayer can provide data but cannot forge inclusion under a trusted finalized source block root.
The destination bridge still trusts the source-chain finality rule, source bridge contract, proof verifier, token mapping, and replay protection.
\end{lstlisting}

This is better, but not magic.

Possible bugs:

\begin{itemize}
\tightlist
\item
  The verifier accepts an unfinalized source block.
\item
  The proof checks receipt inclusion but not emitting contract.
\item
  The event omits destination chain ID.
\item
  The nonce is global instead of scoped by source chain and source
  bridge.
\item
  The token mapping can be changed by an admin without delay.
\item
  The same event can be replayed to another destination bridge.
\item
  The source bridge is upgradeable and can emit fake deposits.
\item
  The proof parser accepts malformed trie nodes.
\item
  The verifier runs out of gas on valid proofs during congestion.
\end{itemize}

The cryptographic proof is only one part of the security model. The
transition after verification is where value moves.

\section{Practical Discipline: Root, Proof, Claim}

When someone says ``we have a proof,'' write three lines:

\begin{lstlisting}
Root:
  What compact commitment is trusted?

Proof:
  What evidence connects the claim to the root?

Claim:
  What exact statement follows if verification succeeds?
\end{lstlisting}

Then write three more:

\begin{lstlisting}
Root trust:
  Why is the root accepted?

Domain:
  Where is the proof valid?

Action:
  What state transition occurs because of the proof?
\end{lstlisting}

This prevents the most common proof-system mistake: verifying a true
fact and then doing an unjustified thing with it.

\section{Review Questions}

\begin{enumerate}
\def\labelenumi{\arabic{enumi}.}
\tightlist
\item
  What does a cryptographic commitment bind, and what does it not prove by itself?
\item
  Why does a proof need a trusted root and a precise verification algorithm?
\item
  What does a Merkle inclusion proof prove relative to a root?
\item
  Why does an inclusion proof not prove finality, semantic correctness, or non-replay by itself?
\item
  What is domain separation, and why does it matter for commitments and proofs?
\item
  How can stale roots, wrong roots, or wrong verification rules create valid-looking but unsafe claims?
\item
  Why are proof verification functions also state-transition surfaces?
\item
  What must a verifier check before accepting a cross-domain proof?
\end{enumerate}

\section{Applied Exercises}

\subsection{Exercise 1: Manual Merkle Verification}

Given four leaves and a claimed proof for one leaf, draw the tree and recompute the root.

Expected output:

\begin{itemize}
\tightlist
\item
  tree diagram;
\item
  proof path;
\item
  recomputation steps;
\item
  final accept/reject decision.
\end{itemize}

\subsection{Exercise 2: Root and Domain Review}

A bridge contract accepts a proof that a deposit event appeared on a source chain. The proof verifies against a block root, but the contract does not check source chain ID, finality depth, message nonce, or whether the message was already consumed.

Identify the failure modes.

Expected output: one table with columns for missing check, false assumption, possible impact, and mitigation.

\subsection{Exercise 3: Proof-Verification Pseudocode}

Write pseudocode for a verifier with this shape:

\begin{lstlisting}
verify(root, claim, proof, domain, nonce) -> accept or reject
\end{lstlisting}

It must check:

\begin{itemize}
\tightlist
\item
  root source;
\item
  domain;
\item
  proof algorithm;
\item
  claim encoding;
\item
  replay or message consumption;
\item
  finality or settlement assumption.
\end{itemize}

Expected output: pseudocode and a short explanation of each rejection condition.

\subsection{Exercise 4: Trust-Minimized Claim Statement}

Rewrite this weak statement:

\begin{lstlisting}
The bridge is trustless because it verifies proofs.
\end{lstlisting}

Write a precise security claim that names the trusted root, proof algorithm, finality assumption, replay protection, and verifier authority.

Expected output: one improved security claim and one residual-risk paragraph.

\section{Key Takeaways}

\begin{itemize}
\tightlist
\item
  Hashes are compact fingerprints of bytes, not proofs of meaning,
  authority, freshness, or finality.
\item
  Preimage, second-preimage, and collision resistance are different
  properties used for different security goals.
\item
  Commitments bind data only if the primitive, encoding, domain
  separation, and verification rules are correct.
\item
  Hash commitments do not hide low-entropy values unless they include
  sufficient secret randomness.
\item
  Merkle proofs prove membership relative to a specific root and tree
  construction.
\item
  State roots, transaction roots, and receipt roots are commitments to
  different parts of blockchain data.
\item
  Ethereum Merkle Patricia proofs and SSZ proofs are different proof
  systems with different verification rules.
\item
  A valid proof against an untrusted, stale, or wrong-domain root is
  useless or dangerous.
\item
  Proof verification is itself a state transition when it gates minting,
  withdrawal, governance, or settlement.
\item
  Trust-minimized does not mean trust-free. It means the remaining
  assumptions are narrower and should be named precisely.
\end{itemize}

\section{Further Reading}

\begin{itemize}
\tightlist
\item
  NIST,
  \href{https://nvlpubs.nist.gov/nistpubs/FIPS/NIST.FIPS.180-4.pdf}{``FIPS
  180-4: Secure Hash Standard''}. Specification of SHA-family hash
  algorithms and their message-digest role.
\item
  NIST,
  \href{https://nvlpubs.nist.gov/nistpubs/Legacy/SP/nistspecialpublication800-107r1.pdf}{``SP
  800-107 Revision 1: Recommendation for Applications Using Approved
  Hash Algorithms''}. Useful legacy definitions of collision, preimage,
  and second-preimage resistance and security-strength guidance; NIST
  notes that supplementary material has moved to its hash-functions
  webpage and that remaining requirements are to move into CMVP
  implementation guidance before withdrawal.
\item
  Satoshi Nakamoto, \href{https://bitcoin.org/bitcoin.pdf}{``Bitcoin: A
  Peer-to-Peer Electronic Cash System''}. See the simplified payment
  verification discussion for Merkle branches and block headers.
\item
  IETF, \href{https://datatracker.ietf.org/doc/html/rfc6962}{``RFC 6962:
  Certificate Transparency''}. A clear standard reference for Merkle
  hash trees, inclusion proofs, and consistency proofs.
\item
  Gavin Wood,
  \href{https://ethereum.github.io/yellowpaper/paper.pdf}{``Ethereum: A
  Secure Decentralised Generalised Transaction Ledger''}. Formal
  Ethereum reference for state, transactions, receipts, and trie roots.
\item
  ethereum.org,
  \href{https://ethereum.org/developers/docs/data-structures-and-encoding/patricia-merkle-trie/}{``Merkle
  Patricia Trie''}. Practical explanation of Ethereum's modified
  Merkle-Patricia trie and state proofs.
\item
  ethereum.org,
  \href{https://ethereum.org/developers/docs/blocks/}{``Blocks''}.
  Useful overview of block fields including
  \passthrough{\lstinline!state\_root!},
  \passthrough{\lstinline!transactionsRoot!}, and
  \passthrough{\lstinline!receiptsRoot!}.
\item
  Ethereum Improvement Proposals,
  \href{https://eips.ethereum.org/EIPS/eip-1186}{``EIP-1186: RPC-Method
  to get Merkle Proofs - eth\_getProof''}. Specification for retrieving
  account and storage proofs.
\item
  Ethereum Consensus Specs,
  \href{https://ethereum.github.io/consensus-specs/ssz/simple-serialize/}{``SimpleSerialize
  (SSZ)''}. Consensus-layer serialization and merkleization reference.
\end{itemize}
